\documentclass[english,11pt,a4paper]{article}
\usepackage[margin=1in]{geometry}
\usepackage{fontspec}
\setmainfont{lmroman10-regular.otf}[BoldFont=lmroman10-bold.otf,ItalicFont=lmroman10-italic.otf,BoldItalicFont=lmroman10-bolditalic.otf]
\setsansfont{lmsans10-regular.otf}[BoldFont=lmsans10-bold.otf]
\setmonofont{lmmono10-regular.otf}[Scale=MatchLowercase]
\newlength{\LeanCharWidth}
\usepackage{amsmath,amssymb}
\usepackage[english]{babel}
\usepackage{microtype,parskip}
\usepackage{graphicx,float,caption}
\usepackage{longtable,booktabs,array,calc}
\usepackage{xcolor}
\usepackage{xurl}
\usepackage{hyperref,bookmark}
\usepackage{etoolbox,needspace}
\setcounter{secnumdepth}{-\maxdimen}
\setlength{\emergencystretch}{3em}
\providecommand{\tightlist}{\setlength{\itemsep}{0pt}\setlength{\parskip}{0pt}}
\newcounter{none}
\AtBeginEnvironment{longtable}{\Needspace{8\baselineskip}\small}
\pretocmd{\section}{\Needspace{5\baselineskip}}{}{}
\pretocmd{\subsection}{\Needspace{4\baselineskip}}{}{}
\pretocmd{\subsubsection}{\Needspace{4\baselineskip}}{}{}
\captionsetup{font=small,labelfont=bf}
\urlstyle{same}
\AtBeginDocument{\special{pdf:minorversion 7}}

\hypersetup{pdflang=en-US,colorlinks=true,urlcolor=blue,linkcolor=blue,pdftitle={Formalization companion for Recovering Host-Triangle Decompositions after Random Deletions in Split Graphs},pdfauthor={Juan Pablo Traverso Gianini}}
\title{Formalization companion for Recovering Host-Triangle Decompositions after Random Deletions in Split Graphs}
\author{Juan Pablo Traverso Gianini\\\small Independent researcher, Santiago, Chile\\\small\href{mailto:jtraverso@gmail.com}{jtraverso@gmail.com}\quad\href{https://orcid.org/0009-0003-6068-4096}{ORCID 0009-0003-6068-4096}}
\date{Preprint 1.0 \enspace·\enspace November 24, 2026}
\begin{document}\raggedbottom\maketitle
\begin{center}Paper III in the series on clique partition gaps in split graphs.\end{center}
\noindent DOI \url{https://doi.org/10.5281/zenodo.24494747}\\
Series repository \url{https://github.com/jtraverso/split-clique-partition-gap}
\subsection{Reading the formal statements}\label{reading-the-formal-statements}

The English preprint uses the source archive dated November 24, 2026 that accompanied drafts 0.4 and 0.5. Section 7 gives its SHA-256 and the main declaration names. The supplementary ZIP is also embedded in the manuscript PDF. The audit driver prints axiom dependencies for 77 declarations in the namespace \texttt{SplitGapIII}.

The selected-support theorem is Theorem 1 in the preprint (Theorem 5 in the supplement's original numbering). The lemma \nolinkurl{hosted_clique_deficit_mem} converts integer zero-margin corrections into a statement about the physical triangle lattice under its switch-availability hypotheses. The lemma \nolinkurl{selected_hinge_condition} turns counts of ordered completions into the edge-set condition H4. The concentration estimates and cut argument producing a single suitable random support appear in §5 and Appendix B.

Theorem 2 (formerly Theorem 6) compares statistics on two scales. The lemma \nolinkurl{hosted_pasch_one_block} identifies a block from the four hosted triangles and the ownership conditions. Given the selected seed, the physical hypotheses and the bounds comparing seed and reference targets, \nolinkurl{selected_decomposition_with_reference_targets} yields a common decomposition. Section 6 establishes the target comparison in the random construction and derives both asymptotic coefficients. Theorem 3 (formerly Theorem 1) is the gap application. The all-orders development in Lean uses extra host copies. Section 2 proves the sharp threshold for a single pool analytically.

\subsection{Reproducing the declaration audit}\label{reproducing-the-declaration-audit}

Extract the source ZIP into an empty directory. The toolchain file specifies Lean 4.35.0-rc3, and the Lake manifest pins Mathlib and the physical absorption library. From the extracted root run

\begin{verbatim}
lake exe cache get
lake build SplitGapIII
lake env lean SplitGapIIIAudit.lean
\end{verbatim}

Compare the declaration names and their axiom lists with \texttt{AUDIT.txt} in the ZIP. Line endings may differ. The preprint validation directory includes a replay against previously compiled dependencies. The full archive checksum also identifies the exact sources and the historical record.

\subsection{Edition concordance}\label{edition-concordance}

{\def\LTcaptype{none} % do not increment counter
\begin{longtable}[]{@{}ll@{}}
\toprule\noalign{}
Review draft 0.5 & Preprint 1.0 \\
\midrule\noalign{}
\endhead
\bottomrule\noalign{}
\endlastfoot
Theorem 5, balanced support & Theorem 1 \\
Theorem 6, statistical recovery & Theorem 2 \\
Theorem 1, gap and connectivity & Theorem 3 \\
Theorem 2, random host lists & Theorem 4 \\
Lemma 3, two coordinate orders & Lemma 5 \\
Proposition 4, deterministic budgets & Proposition 6 \\
Lemma 6.1, localization of Pasch configurations & Lemma 7 \\
\end{longtable}
}

All thirteen numbered displays retain their draft numbering. The manuscript and the source archive are available in the series repository. \url{https://github.com/jtraverso/split-clique-partition-gap/tree/main/preprints/SPLITGAPIII}

\end{document}
